\chapter{Superseded Solver-Led Words Draft}
\label{ch:words}

The byte \code{11111111} is eight marks on a page. A processor can hold the
same pattern in eight switches. We may read it as the natural number 255, the
signed integer -1, a mask with every switch on, or the result of an
exceptional operation. The bits remain unchanged while their mathematical
role changes. They do not choose their own meaning. An operation and its
specification do.

Begin with a four-bit counter. It visits 0 through 15. Add one to 15, and the
counter returns to 0. There is no fifth position to hold the 16. Add two, and
it stops at 1. The machine has not made an arithmetic mistake. It is counting
in a finite circle.

That small circle is the first mathematical object we need.

\section{The word of width \(n\)}
\label{sec:word}

\begin{definition}[title={Word \(n\)}]
Let \(n \geq 1\). The \keyterm{word of width \(n\)} is the set of residues
modulo \(2^n\), written \(\mathbf{Z}/2^n\mathbf{Z}\). We represent its
elements by the natural numbers \(0,1,\ldots,2^n-1\). Addition and
multiplication use the ordinary result and then keep its remainder modulo
\(2^n\).

For \(0 \leq i < n\), bit \(i\) of a word \(x\) is
\((x \div 2^i) \bmod 2\). Bitwise \emph{and}, \emph{or}, and \emph{xor}
apply their Boolean operation at each bit position. A left shift by \(k\)
multiplies by \(2^k\) and reduces modulo \(2^n\). A logical right shift by
\(k\) takes the quotient after division by \(2^k\), rounded down to a whole
number.
\end{definition}

Here \(\mathbf{Z}\) means the integers. The quotient notation groups together
integers that differ by a multiple of \(2^n\). No special notation is needed
to calculate. Take the ordinary result and keep the remainder after division
by \(2^n\).

At width four, this rule turns \(15+1\) into 0 and \(15+2\) into 1. It does
not pretend that the natural numbers 16 and 17 have vanished. It says that 16
and 0 represent the same stored word at this width, as do 17 and 1.

The construction is small, but its consequences spread through the book. A
word of width \(n\) has exactly \(2^n\) possible values. An instruction on one
register is therefore a function on a finite set, or on a product of finite
sets when it takes several inputs. Every later statement about instruction
equivalence compares such functions. Every later exhaustive search depends
on the fact that their domains have edges we can name and, sometimes, count.

There is beauty in this compression. The same \(n\) bits can be printed as a
string, stored as a physical state, read as a number, and passed as the input
to a function. Those descriptions are not identical, but an exact semantics
lets us move among them without losing the object.

Axeyum is the reasoning software used to check this book's solver-backed
objects. It constructs the natural numbers used here without assumptions. The
present stack can build fixed-width bit-vector terms from those naturals, but
it does not yet package this definition as one reusable word library inside
the kernel. This limit matters: rebuilding fixed widths is not the same as
stating and checking the laws once for every positive width.

\begin{artifact}{W.def.word}
This object fixes the mathematical carrier used throughout the book. The
current kernel route rebuilds a word at the needed width. A reusable word
prelude, with the operations and laws developed once for every positive width,
is tracked by \obj{OP.kernel.word-prelude}.
\ArtifactScope{W.def.word}
\end{artifact}

We now know which values fit in a word. We do not yet know what every
operation should do at its awkward edges. Division gives us the first choice.

\section{Why division needs an answer}
\label{sec:totality}

In ordinary integer arithmetic, \(7\div2\) has a quotient and a remainder.
The expression \(7\div0\) has neither. A machine specification cannot leave
the corresponding bit-vector operation half described and still expect a
compiler, solver, and hardware model to compare the same function.

Fixed-size bit-vector logic therefore makes its operations total: every pair
of input words receives an output word. Totality does not force one particular
answer to division by zero. It forces the specification to choose an answer.
Once chosen, the awkward edge becomes as exact as any ordinary input.

SMT-LIB's fixed-size bit-vector theory chooses two base rules. Unsigned
division by zero returns the all-ones word. Unsigned remainder by zero returns
the dividend. Signed division is defined from unsigned division and
two's-complement signs. With a zero divisor, it returns 1 for a negative
dividend and the all-ones word otherwise \autocite{smtlibfixedbv}.

At width four, all ones is \code{1111}. Its unsigned reading is 15. Under the
signed two's-complement reading it is -1: because the high bit is one, subtract
\(2^4\) from 15. One pattern now serves two statements. Unsigned division by
zero returns 15; the zero-or-positive branch of signed division by zero returns
-1. The returned bits agree. The readings differ.

These rules are design conventions, not consequences of ordinary division.
Yet after we adopt the standard as our semantics, they become mathematical
claims about a completely specified function. We can ask whether they hold
for one example, for every byte, or for every positive width. Those are three
different questions.

\section{Turning every input into one question}
\label{sec:refuting-opposite}

For a byte \(x\), the unsigned claim is

\[
x\div_u0=255
\qquad\mbox{for every }x\in\mathbf{Z}/256\mathbf{Z}.
\]

The reader proof is one line once the operation has been defined. The divisor
is zero, so the zero-divisor branch returns the word whose eight bits are all
one. Under the unsigned reading, that word is 255. The input \(x\) never
affects the selected branch or its result.

The machine route approaches the same sentence from the other side. It asks
for a \keyterm{counterexample}: one byte \(x\) such that
\(x\div_u0\ne255\). If the solver finds one, the original claim fails. If the
encoded formula has no solution, no byte in the width-eight domain violates
the claim. The universal statement has become a search for one exception.

\begin{artifact}{W.thm.udiv0}
Its artifact, \code{bvudiv-by-zero-is-allones.smt2}, contains the negated
width-eight statement. The checker sends it through Axeyum's bit-vector solver
route and expects the verdict that no counterexample exists.
\ArtifactScope{W.thm.udiv0}
\end{artifact}

The two proofs meet at the zero-divisor branch. A reader can see why the input
cannot matter. The machine route checks every byte under the actual
encoding used by the book. Neither tells us why the standard's designers
preferred this result. Mathematics can expose the consequence of a convention
without inventing a necessity behind the choice.

This modest example contains a pattern of surprising reach. To establish a
statement about 256 inputs, we do not recite 256 successful calculations. We
describe one function, turn failure into a formula, and ask whether any
counterexample exists. The finite domain remains real; the proof makes its
structure do the work.

\section{Two neighboring rules}
\label{sec:neighboring-rules}

Unsigned remainder uses a different zero-divisor branch. At width four,
\(x\bmod_u0=x\). The reader proof again selects the defining branch, but this
time the returned word depends on the dividend. The width-eight machine query
asks whether any byte differs from its own remainder by zero.

\begin{artifact}{W.thm.urem0}
This artifact asks for a width-eight counterexample to \(x\bmod_u0=x\). The
same Axeyum route checks the formula.
\ArtifactScope{W.thm.urem0}
\end{artifact}

Signed division adds a case split. Read the high bit as the sign. If it is
one, the dividend is negative and division by zero returns the word for 1. If
it is zero, the operation returns all ones, whose signed reading is -1. Every
four-bit word enters exactly one branch. That split is the reader proof. At
width eight, the formula represents the same split with a signed comparison.

\begin{artifact}{W.thm.sdiv0}
This artifact asks for a byte on which the signed zero-divisor case differs
from the two-branch result. Its scope names both the width and the SMT-LIB
semantics being checked.
\ArtifactScope{W.thm.sdiv0}
\end{artifact}

All three arguments have the same outer shape: choose the zero-divisor branch,
state its result, and search for a counterexample. Inside that shape, the
functions differ. One returns all ones, one returns its input, and one chooses
between two results by a sign. The phrase \emph{division by zero} is too vague
to preserve these distinctions. The exact operation makes them visible.

\section{Make the checker lose}
\label{sec:negative-control}

Now give the same route a false rule: unsigned division by zero returns zero.
At width four, the specified branch returns 15 instead. Every dividend exposes
the error. At width eight, the solver should find a counterexample rather than
report that none exists.

A \keyterm{negative control} is this deliberate failure. It tests whether the
checking route distinguishes the expected claim from a nearby false one. The
control is evidence about the checker, not another proof of the theorem.

\begin{keyidea}[title={A checker that cannot fail}]
If a checker reports the desired answer for every input, a true claim cannot
distinguish it from a broken checker. A false neighbor can. The book therefore
sends positive evidence and a negative control through the same route.
\end{keyidea}

\begin{artifact}{W.ctl.udiv0-wrong}
This object asserts the negation of the false zero-result rule. The checker
expects a counterexample to exist; reversing that expectation is itself the
control that must fail.
\ArtifactScope{W.ctl.udiv0-wrong}
\end{artifact}

This failed control has its own austere elegance. We learn something about a
successful check by arranging a neighboring check that must not succeed.
Correct rejection becomes evidence. The machinery earns trust not by being
infallible, but by showing one precise way in which it can say no.

That attack still has a boundary. It catches a route that gives the same
verdict to every formula. It cannot catch every shared mistake. If the theorem
and its control use the same bad encoding of division, both may behave as
expected for the wrong reason. Independent semantics or a second checker
would narrow that trust boundary.

\section{What the verdict carries}
\label{sec:verdict}

All three theorem objects use a \keyterm{front-door verdict}. The command-line
program reads an SMT-LIB formula and reports whether a counterexample exists.
Axeyum checks the solver result inside that route, but the command used for
these objects does not export a separate certificate for another checker to
read. The reproducible verdict and firing control are the evidence we have.
An exported-certificate claim would require another object and another route.

The boundary has three layers. These formulas cover width eight. Axeyum's
current front door checks them but exports no certificate. Its kernel does not
yet contain a reusable word theory that states these laws once for every
positive width. None of those limits weakens the width-eight statement. Each
prevents us from silently making it stronger.

We began with eight marks that did not choose their own meaning. We end with a
specified function, a statement about every byte, two ways to support that
statement, and a deliberate false neighbor that the checker rejects. The
distance between those points is the work of semantics and proof.

One word is not yet a machine. A running program also needs places for words,
a place in the program, memory, and a record of exceptional outcomes. Chapter
2 assembles those parts into a state that an instruction can change.

\section{Exercises}

\begin{enumerate}
\item List all sixteen width-four words. Check by hand that unsigned division
by zero returns \code{1111} for each one, then reproduce the width-eight
object \obj{W.thm.udiv0}.
\item Change the expected result in the unsigned-division formula from all
ones to zero. Predict a counterexample before running the checker. Explain why
the changed formula is a useful control but not evidence for the original
claim by itself.
\item State \obj{W.thm.udiv0} for every positive width at once. List the pieces
a reusable kernel proof would need that 256 fixed-width cases do not supply.
\item Show that there are \(2^n\) words of width \(n\) and
\(2^{n2^n}\) functions from words to words. Explain why this large function
count does not make the division-by-zero proof difficult.
\item Design a different total result for unsigned division by zero. Name one
definition, one compiler transformation, and one checker expectation that
would have to change so that all three still describe the same operation.
\item The same eight bits can represent 255, -1, a mask, or an exceptional
result. Give one additional interpretation. State what information must be
supplied before the bit pattern acquires that meaning.
\end{enumerate}
